% !TeX program = lualatex
% halo:begin
% {
%   "version": 1,
%   "publish": true,
%   "course": "mathematical-analysis-i",
%   "section": "2.1",
%   "title": "数学分析（一）2.1：数列极限",
%   "categories": ["study-notes"],
%   "tags": ["mathematics"],
%   "excerpt": "数列极限的定义、无穷小量与无穷大量、基本例题，以及习题 2.1 的课后作业。"
% }
% halo:end
% 来源：IMG_1262.HEIC、IMG_1263.HEIC、IMG_1264.HEIC；
% 本次补充依据 IMG_1365.HEIC、IMG_1366.HEIC、IMG_1367.HEIC、
% IMG_1369.HEIC 上半页，作业题干对照《数学分析》第五版上册习题 2.1。
\RequirePackage{iftex}
\RequireLuaTeX
\documentclass[UTF8, fontset=fandol, a4paper, 12pt]{ctexart}
\usepackage[margin=2.5cm]{geometry}
\usepackage{amsmath,amssymb,amsthm}
\usepackage{enumitem,fancyhdr}
\usepackage[hidelinks]{hyperref}
\newcommand{\R}{\mathbb R}
\newcommand{\Nplus}{\mathbb N_{+}}
\theoremstyle{definition}
\newtheorem{definition}{定义}[section]
\newtheorem{example}{例题}
\renewcommand{\proofname}{证明}
\setlist[enumerate]{leftmargin=2em,itemsep=0.2em,topsep=0.3em}
\setlength{\headheight}{16pt}
\pagestyle{fancy}
\fancyhf{}
\fancyhead[L]{\small 数学分析（一）}
\fancyhead[R]{\small 2.1\quad 数列极限}
\fancyfoot[C]{\thepage}
\renewcommand{\headrulewidth}{0.3pt}
\raggedbottom
\title{数列极限}
\author{Yuchen Zhang}
\date{整理日期：2026年9月21日}
\makeatletter
\renewcommand{\maketitle}{\begin{center}
{\LARGE\bfseries \@title\par}\vspace{0.6em}
\@author\par\vspace{0.2em}{\small\@date\par}
\end{center}}
\makeatother
\begin{document}
\maketitle
\section{数列极限的定义}
\begin{definition}[收敛与极限]
设 $\{a_n\}$ 为实数列，$A\in\R$。若
\[
 \forall\varepsilon>0,\quad\exists N\in\Nplus,\quad
 \forall n>N,\quad |a_n-A|<\varepsilon,
\]
则称数列 $\{a_n\}$ 收敛于 $A$，记为
\[
 \lim_{n\to\infty}a_n=A\qquad\text{或}\qquad a_n\to A\ (n\to\infty).
\]
\end{definition}
\begin{enumerate}
\item $\varepsilon$ 可以任意小，但必须为正；$N$ 可以依赖于 $\varepsilon$，
也可以依赖于给定数列与 $A$，但不能依赖于随后选取的 $n$。
\item $N$ 不唯一：某个 $N$ 可用时，任何更大的正整数也可用。
通常先由 $|a_n-A|<\varepsilon$ 估计出对 $n$ 的要求，再选取整数 $N$。
\item 对固定的 $c>0$、$p>0$，定义右端的 $\varepsilon$ 可以换成
$c\varepsilon$ 或 $\varepsilon^p$，也可将严格不等号改为 $\le$，所得定义均等价。
例如，若能使误差小于 $c\varepsilon$，则对目标误差 $\eta>0$ 取
$\varepsilon=\eta/c$ 即可；若误差仅保证 $\le\varepsilon$，则取 $\varepsilon=\eta/2$。
\end{enumerate}

\subsection{邻域刻画}
记 $U(A;\varepsilon)=(A-\varepsilon,A+\varepsilon)$，则
\[
 |a_n-A|<\varepsilon\iff A-\varepsilon<a_n<A+\varepsilon
 \iff a_n\in U(A;\varepsilon).
\]
因此，$a_n\to A$ 当且仅当：对任意 $\varepsilon>0$，落在
$U(A;\varepsilon)$ 外的数列项只有有限项。
这里按下标计数，即集合 $\{n\in\Nplus:a_n\notin U(A;\varepsilon)\}$ 有限。
一方面，$n>N$ 时均在邻域内，邻域外至多有前 $N$ 项；另一方面，
若邻域外仅有有限项，取 $N$ 不小于这些项的所有下标即可（没有这样的项时取 $N=1$）。

\subsection{发散及其量词表述}
数列不收敛于任何实数时，称其发散。准确地说，
\[
 \forall A\in\R,\quad\exists\varepsilon_0>0,\quad
 \forall N\in\Nplus,\quad\exists n>N:\ |a_n-A|\ge\varepsilon_0.
\]
即对每个候选极限 $A$，都存在一个邻域，其外有无穷多项。
只排除某一个 $A$，不能证明数列发散。

\newpage
\section{由定义证明极限：基本例题}
\begin{example}
证明 $\displaystyle\lim_{n\to\infty}\frac1n=0$。
\end{example}
\begin{proof}
任取 $\varepsilon>0$，取 $N=\max\{1,\lceil1/\varepsilon\rceil\}$。
对任意 $n>N$，有
\[
 \left|\frac1n-0\right|=\frac1n<\frac1N\le\varepsilon.
\]
故结论成立。这里 $\lceil t\rceil$ 为不小于 $t$ 的最小整数。
\end{proof}

\begin{example}
证明 $\displaystyle\lim_{n\to\infty}\frac{3n^2}{n^2-3}=3$。
\end{example}
\begin{proof}
当 $n\ge2$ 时，
\[
 \left|\frac{3n^2}{n^2-3}-3\right|=\frac9{n^2-3}.
\]
任取 $\varepsilon>0$，取
\[
 N=\max\left\{2,\left\lceil\sqrt{\frac9\varepsilon+3}\right\rceil\right\}.
\]
对任意 $n>N$，有 $n^2>9/\varepsilon+3$，于是
\[
 \left|\frac{3n^2}{n^2-3}-3\right|=\frac9{n^2-3}<\varepsilon.
\]
故结论成立。
\end{proof}

\begin{example}
设 $q\in\R$ 且 $|q|<1$，证明 $\displaystyle\lim_{n\to\infty}q^n=0$。
\end{example}
\begin{proof}
$q=0$ 时显然成立。以下设 $0<|q|<1$。
任取 $\varepsilon>0$，因 $\ln|q|<0$，有
\[
 |q|^n<\varepsilon
 \iff n\ln|q|<\ln\varepsilon
 \iff n>\frac{\ln\varepsilon}{\ln|q|}.
\]
取
\[
 N=\max\left\{1,\left\lceil\frac{\ln\varepsilon}{\ln|q|}\right\rceil\right\}.
\]
则对任意 $n>N$，都有 $|q^n-0|=|q|^n<\varepsilon$，故结论成立。
取最大值保证了 $\varepsilon\ge1$ 时 $N$ 仍是正整数。
\end{proof}

\newpage
\section{根式与方根的极限}
\begin{example}
设 $x_n\to a>0$，证明 $\sqrt{x_n}$ 从充分大的 $n$ 起有定义，且
$\sqrt{x_n}\to\sqrt a$。
\end{example}
\begin{proof}
先保证根式有意义。由极限定义，存在 $N_0\in\Nplus$，使 $n>N_0$ 时
$|x_n-a|<a/2$，从而 $x_n>a/2>0$。
因此 $\sqrt{x_n}$ 至少从某一项起有定义；若要把它看成从第一项起的实数列，
可要求所有 $x_n\ge0$，或将前面有限项另行赋值，均不影响极限。

任取 $\varepsilon>0$，由 $x_n\to a$，存在 $N_1\in\Nplus$，
使 $n>N_1$ 时 $|x_n-a|<\varepsilon\sqrt a$。
取 $N=\max\{N_0,N_1\}$，则对任意 $n>N$，
\[
 |\sqrt{x_n}-\sqrt a|
 =\frac{|x_n-a|}{\sqrt{x_n}+\sqrt a}
 \le\frac{|x_n-a|}{\sqrt a}<\varepsilon.
\]
故 $\sqrt{x_n}\to\sqrt a$。
\end{proof}
\textbf{另一种估计。} 对 $u,v\ge0$，由
$|u-v|=|\sqrt u-\sqrt v|(\sqrt u+\sqrt v)$ 及
$|\sqrt u-\sqrt v|\le\sqrt u+\sqrt v$，可得
\[
 |\sqrt u-\sqrt v|\le\sqrt{|u-v|}.
\]
故只要使 $|x_n-a|<\varepsilon^2$，就能使
$|\sqrt{x_n}-\sqrt a|<\varepsilon$。

\begin{example}
设 $a>0$，证明 $\displaystyle\lim_{n\to\infty}a^{1/n}=1$。
\end{example}
\begin{proof}
$a=1$ 时显然成立。
若 $a>1$，令 $t_n=a^{1/n}-1>0$。由二项式展开，
\[
 a=(1+t_n)^n\ge1+nt_n,\qquad 0<t_n\le\frac{a-1}{n}.
\]
任取 $\varepsilon>0$，取 $N=\max\{1,\lceil(a-1)/\varepsilon\rceil\}$，
则 $n>N$ 时，$|a^{1/n}-1|\le(a-1)/n<\varepsilon$。

若 $0<a<1$，令 $b=1/a>1$。由前面的估计，
\[
 |a^{1/n}-1|=1-\frac1{b^{1/n}}
 =\frac{b^{1/n}-1}{b^{1/n}}\le b^{1/n}-1\le\frac{b-1}{n}.
\]
取 $N=\max\{1,\lceil(b-1)/\varepsilon\rceil\}$ 即得所需结论。
\end{proof}

\newpage
\subsection*{补充例题：$\sqrt[n]{n}$ 的极限}
证明 $\displaystyle\lim_{n\to\infty}n^{1/n}=1$。
\begin{proof}
令 $t_n=n^{1/n}-1\ge0$。当 $n\ge2$ 时，二项式展开给出
\[
 n=(1+t_n)^n\ge\binom n2t_n^2,
 \qquad 0\le t_n\le\sqrt{\frac2{n-1}}.
\]
任取 $\varepsilon>0$，取 $N=\max\{2,\lceil1+2/\varepsilon^2\rceil\}$。
对任意 $n>N$，有
\[
 |n^{1/n}-1|=t_n\le\sqrt{\frac2{n-1}}<\varepsilon.
\]
所以 $n^{1/n}\to1$。此处底数随 $n$ 变化，不能直接套用固定 $a$ 的结论。
\end{proof}

\section{阶乘与幂的比较}
\begin{example}
证明 $\displaystyle\lim_{n\to\infty}\frac{2^n}{n!}=0$；
进一步，对任意固定的 $\alpha\in\R$，证明
$\displaystyle\lim_{n\to\infty}\frac{\alpha^n}{n!}=0$。
\end{example}
\begin{proof}
先看 $2^n/n!$。当 $n\ge2$ 时，
\[
 \frac{2^n}{n!}
 =\frac4n\prod_{k=2}^{n-1}\frac2k\le\frac4n,
\]
其中 $n=2$ 时空乘积取 $1$。
任取 $\varepsilon>0$，取 $N=\max\{2,\lceil4/\varepsilon\rceil\}$，
则 $n>N$ 时，$|2^n/n!|\le4/n<\varepsilon$。

一般地，$\alpha=0$ 时显然成立。若 $\alpha\ne0$，选取整数
$m\ge\max\{1,2|\alpha|\}$，并记 $C=|\alpha|^m/m!>0$。
对 $n>m$，有
\[
 \left|\frac{\alpha^n}{n!}\right|
 =C\prod_{k=m+1}^{n}\frac{|\alpha|}{k}
 \le C\left(\frac12\right)^{n-m}.
\]
给定 $\varepsilon>0$，取
\[
 N=m+\max\{1,\lceil\log_2(C/\varepsilon)\rceil\}.
\]
对任意 $n>N$，有 $n-m>\log_2(C/\varepsilon)$，于是
$|\alpha^n/n!|\le C2^{-(n-m)}<\varepsilon$，结论成立。
\end{proof}

\newpage
\section{用定义证明数列发散}
\begin{example}
证明数列 $\{n^2\}$ 与 $\{(-1)^n\}$ 均发散。
\end{example}
\begin{proof}
\textbf{（1）$\{n^2\}$。}
任取 $A\in\R$，取 $\varepsilon_0=1$。
对于任意 $N\in\Nplus$，选取正整数
\[
 n>\max\{N,1,A+1\}.
\]
则 $n^2\ge n>A+1$，从而
\[
 |n^2-A|=n^2-A>1=\varepsilon_0.
\]
所以 $A$ 不是该数列的极限。由于 $A$ 任意，$\{n^2\}$ 发散。

\textbf{（2）$\{(-1)^n\}$。}
任取 $A\in\R$，仍取 $\varepsilon_0=1$。
\begin{itemize}[leftmargin=2em]
\item 若 $A\ge0$，则对任意 $N$，可取奇数 $n>N$，此时
\[
 |(-1)^n-A|=|-1-A|=1+A\ge1.
\]
\item 若 $A<0$，则对任意 $N$，可取偶数 $n>N$，此时
\[
 |(-1)^n-A|=|1-A|=1-A>1.
\]
\end{itemize}
所以对每个 $A\in\R$，邻域 $U(A;1)$ 外都有无穷多项，故该数列发散。
\end{proof}

\newpage
\section{无穷小量与无穷大量}
\begin{definition}[无穷小量]
极限为 $0$ 的数列称为无穷小量。数列 $a_n\to a$ 当且仅当
$\alpha_n=a_n-a$ 是无穷小量；此时可写成 $a_n=a+\alpha_n$。
\end{definition}

\begin{definition}[无穷大量]
若数列 $\{a_n\}$ 满足
\[
 \forall M>0,\quad\exists N\in\Nplus,\quad
 \forall n>N,\quad |a_n|>M,
\]
则称其为无穷大量，记作 $a_n\to\infty$。
若上式中的 $|a_n|>M$ 分别可加强为 $a_n>M$ 或 $a_n<-M$，
则记作 $a_n\to+\infty$ 或 $a_n\to-\infty$。
\end{definition}
记号 $a_n\to\infty$ 只要求绝对值最终超过任意给定的界，
不要求数列最终同号。例如 $(-1)^n n\to\infty$，
但既不趋于 $+\infty$，也不趋于 $-\infty$。

\subsection*{无界与无穷大的区别}
无界只要求对任意 $M>0$，\emph{存在某一项}的绝对值大于 $M$；
无穷大量要求\emph{从某一项起每一项}都如此。
数列 $b_n=[1+(-1)^n]n$ 无界，但其奇数项恒为 $0$，
所以不是无穷大量。无穷大量必发散，因为收敛数列从某项起
落在某个有限区间内，再加上有限个首项即有界。

\newpage
\section{课后作业}
\subsection{习题 2.1，第 1 题}
\noindent\textbf{题目：}设
$a_n=\dfrac{1+(-1)^n}{n}$，候选极限 $a=0$。
（1）对 $\varepsilon_1=0.1$、$\varepsilon_2=0.01$、
$\varepsilon_3=0.001$，分别求出符合极限定义的 $N$；
（2）这是否已经证明 $a_n\to0$；
（3）对任意 $\varepsilon>0$，是否都能找到这样的 $N$？

\begin{proof}
奇数项为 $0$，偶数项为 $2/n$，故 $0\le a_n\le2/n$。
（1）依次可取 $N=20,200,2000$；当 $n>N$ 时，
$|a_n| \le 2/n<2/N=\varepsilon_i$。
（2）仅检查三个指定的 $\varepsilon$，还不足以证明极限，
因为定义要求\emph{每个}正 $\varepsilon$ 都成立。
（3）任给 $\varepsilon>0$，取
$N=\max\{1,\lceil2/\varepsilon\rceil\}$。
对一切 $n>N$，都有 $|a_n|\le2/n<\varepsilon$，因此 $a_n\to0$。
\end{proof}

\subsection{习题 2.1，第 2 题（1）（3）（5）}
\noindent\textbf{题目：}按 $\varepsilon$--$N$ 定义证明：
\[
 \text{（1）}\ \frac{n}{n+1}\to1,\qquad
 \text{（3）}\ \frac{n!}{n^n}\to0,\qquad
 \text{（5）}\ \frac{n}{a^n}\to0\quad(a>1).
\]
\begin{proof}
\textbf{（1）}任取 $\varepsilon>0$，取
$N=\max\{1,\lceil1/\varepsilon\rceil\}$。当 $n>N$ 时，
\[
 \left|\frac{n}{n+1}-1\right|=\frac1{n+1}
 <\frac1n<\varepsilon.
\]
\textbf{（3）}对 $n\ge1$，
\[
 0<\frac{n!}{n^n}
 =\frac1n\cdot\frac2n\cdots\frac nn\le\frac1n.
\]
任取 $\varepsilon>0$，取 $N=\max\{1,\lceil1/\varepsilon\rceil\}$，
则 $n>N$ 时上式小于 $\varepsilon$。

\textbf{（5）}写 $a=1+h$，其中 $h>0$。二项式定理给出
$a^n\ge\binom n2h^2$，所以 $n\ge2$ 时
\[
 0<\frac n{a^n}\le\frac{2}{(n-1)h^2}.
\]
给定 $\varepsilon>0$，取
$N=\max\{2,\lceil1+2/(\varepsilon h^2)\rceil\}$，
则 $n>N$ 时右端小于 $\varepsilon$。
\end{proof}

\subsection{习题 2.1，第 3 题}
\noindent\textbf{题目：}根据本节已证明的基本极限，求
\[
\begin{array}{llll}
\text{（1）}\ \displaystyle\lim\frac1{\sqrt n},&
\text{（2）}\ \displaystyle\lim\sqrt[n]{3},&
\text{（3）}\ \displaystyle\lim\frac1{n^3},&
\text{（4）}\ \displaystyle\lim\frac1{3^n},\\[0.5em]
\text{（5）}\ \displaystyle\lim\frac1{\sqrt{2^n}},&
\text{（6）}\ \displaystyle\lim\sqrt[n]{10},&
\text{（7）}\ \displaystyle\lim\frac1{\sqrt[n]{2}},
\end{array}
\qquad(n\to\infty)
\]
并指出哪些是无穷小数列。

\begin{proof}
（1）对任意 $\varepsilon>0$，当 $n>1/\varepsilon^2$ 时，
$1/\sqrt n<\varepsilon$，故极限为 $0$。
（2）与（6）由前文固定正数的方根极限得到，分别为 $1$ 与 $1$；
更具体地，$0\le\sqrt[n]3-1\le2/n$、
$0\le\sqrt[n]{10}-1\le9/n$。
（3）因 $0<1/n^3\le1/n$，极限为 $0$。
（4）与（5）分别是公比 $1/3$ 与 $1/\sqrt2$ 的几何数列，
两公比的绝对值都小于 $1$，故极限均为 $0$。
（7）令 $b_n=\sqrt[n]2\ge1$，则
\[
 0\le1-\frac1{b_n}=\frac{b_n-1}{b_n}
 \le b_n-1\le\frac1n\to0,
\]
故极限为 $1$。因此（1）、（3）、（4）、（5）是无穷小数列。
\end{proof}

\subsection{习题 2.1，第 4 题}
\noindent\textbf{题目：}证明：若 $a_n\to a$，则对任一正整数 $k$，
$a_{n+k}\to a$。
\begin{proof}
任取 $\varepsilon>0$，由 $a_n\to a$，存在 $N\in\Nplus$，
使 $m>N$ 时 $|a_m-a|<\varepsilon$。
若 $n>N$，则 $n+k>N$，于是
$|a_{n+k}-a|<\varepsilon$。故 $a_{n+k}\to a$。
\end{proof}

\subsection{习题 2.1，第 5 题}
\noindent\textbf{题目：}用“任一候选极限的某个邻域外总有无穷多项”
这一表述证明：（1）$\{1/n\}$ 不以 $1$ 为极限；
（2）$\{n^{(-1)^n}\}$ 发散。
\begin{proof}
\textbf{（1）}取 $\varepsilon_0=1/2$。无论给定哪个 $N$，
都可取 $n>\max\{N,2\}$，使
$|1/n-1|=1-1/n>1/2$。
因此 $1$ 不是 $\{1/n\}$ 的极限。

\textbf{（2）}任取候选极限 $A\in\R$，取 $\varepsilon_0=1$。
对任意 $N$，可取偶数 $n>\max\{N,|A|+1\}$。
此时 $n^{(-1)^n}=n$，且
$|n-A|\ge n-|A|>1$。
所以没有任何实数是该数列的极限，数列发散。
\end{proof}

\subsection{习题 2.1，第 8 题}
\noindent\textbf{题目：}若 $a_n\to a$，证明 $|a_n|\to|a|$；
求使逆命题对任意数列成立的 $a$。
\begin{proof}
由反三角不等式，
\[
 \bigl||a_n|-|a|\bigr|\le|a_n-a|\to0,
\]
故 $|a_n|\to|a|$。若 $a=0$，则
$\bigl||a_n|-|a|\bigr|=|a_n|=|a_n-a|$，逆命题也成立。
若 $a\ne0$，取常数列 $a_n=-a$，有
$|a_n|=|a|$，但 $a_n$ 不趋于 $a$。
因此所求的值恰为 $a=0$。
\end{proof}

\subsection{习题 2.1，第 9 题}
\noindent\textbf{题目：}按 $\varepsilon$--$N$ 定义证明：
\[
 \text{（1）}\ \sqrt{n+1}-\sqrt n\to0,\qquad
 \text{（2）}\ \frac{1+2+\cdots+n}{n^3}\to0,
\]
\[
 \text{（3）}\ a_n\to1,\quad
 a_n=\begin{cases}
 (n-1)/n,&n\text{ 为偶数},\\
 \sqrt{n^2+n}/n,&n\text{ 为奇数}.
 \end{cases}
\]
\begin{proof}
\textbf{（1）}有
\[
 0<\sqrt{n+1}-\sqrt n
 =\frac1{\sqrt{n+1}+\sqrt n}<\frac1{\sqrt n}.
\]
给定 $\varepsilon>0$，取
$N=\max\{1,\lceil1/\varepsilon^2\rceil\}$ 即可。

\textbf{（2）}由等差数列求和式，对 $n\ge1$ 有
\[
 0<\frac{1+2+\cdots+n}{n^3}
 =\frac{n+1}{2n^2}\le\frac1n.
\]
取 $N=\max\{1,\lceil1/\varepsilon\rceil\}$ 即可。

\textbf{（3）}若 $n$ 为偶数，则 $|a_n-1|=1/n$；
若 $n$ 为奇数，则
\[
 |a_n-1|=\sqrt{1+\frac1n}-1
 =\frac{1/n}{\sqrt{1+1/n}+1}\le\frac1{2n}.
\]
故所有 $n$ 都有 $|a_n-1|\le1/n$。
取 $N=\max\{1,\lceil1/\varepsilon\rceil\}$ 即得 $a_n\to1$。
\end{proof}

\subsection{习题 2.1，第 10 题}
\noindent\textbf{题目：}设每个 $a_n\ne0$，证明
$a_n\to0$ 的充要条件是 $1/a_n\to\infty$。
这里“趋于 $\infty$”指绝对值最终超过任意正数，
不预设 $a_n$ 的符号。
\begin{proof}
若 $a_n\to0$，任取 $M>0$，对 $\varepsilon=1/M$ 应用极限定义，
可得充分大 $n$ 时 $|a_n|<1/M$。因为 $a_n\ne0$，
所以 $|1/a_n|>M$，即 $1/a_n\to\infty$。

反之，若 $1/a_n\to\infty$，任取 $\varepsilon>0$，
对 $M=1/\varepsilon$ 应用无穷大量的定义，
可得充分大 $n$ 时 $|1/a_n|>1/\varepsilon$，
从而 $|a_n|<\varepsilon$。故 $a_n\to0$。
\end{proof}
\end{document}
